\documentclass[11pt]{article}

\usepackage[margin=1in]{geometry}
\usepackage{amsmath,amssymb,amsthm}
\usepackage{mathtools}
\usepackage{booktabs}
\usepackage{microtype}
\usepackage[hidelinks]{hyperref}
\usepackage[round]{natbib}
\usepackage{enumitem}
\setlength{\emergencystretch}{2em}

\theoremstyle{plain}
\newtheorem{theorem}{Theorem}
\newtheorem{lemma}{Lemma}
\newtheorem{corollary}{Corollary}
\theoremstyle{definition}
\newtheorem{axiom}{Axiom}
\newtheorem{definition}{Definition}
\theoremstyle{remark}
\newtheorem{remark}{Remark}

\newcommand{\R}{\mathbb{R}}
\newcommand{\E}{\mathbb{E}}
\newcommand{\KL}[2]{D\!\left(#1\,\|\,#2\right)}
\newcommand{\I}{I}
\DeclareMathOperator{\Cov}{Cov}
\DeclareMathOperator*{\argmax}{arg\,max}

\title{\bf A Mathematical Theory of Value}
\author{Cheng Qian}
\date{2026}

\begin{document}
\maketitle

\begin{abstract}
We propose that \emph{value} --- the quantity that goal-directed agents create, destroy, and exchange --- is a
lawful structural quantity, in the same category as information, once stripped of its semantic clothing
(morality, price, psychology). Following the method of \citet{shannon1948}, we make one ruthless abstraction:
value is the rate at which an agent converts a physical resource into goal-progress, relative to a frame fixed
by the agent's goal. A scale-invariance axiom forces a logarithmic measure of value, $V=\sum_i k_i\ln e_i$, via a Cauchy
functional equation; the compounding dynamics of a reinvested resource force the same form independently via
the ergodicity argument of \citet{peters2019}. Two routes, no shared premises --- the over-determination is
the structural signal. From the compounding dynamics we also derive a \emph{coding theorem of value}: the rate at which an agent can create value through a perception channel $Y$ of
the world $X$ is bounded by the mutual information, $\Delta G \le \I(X;Y)$, achieved by Bayes-proportional
allocation; and realized value decomposes exactly as available potential minus dissipation,
$G=\KL{q}{r}-\KL{q}{p}$, identifying misalignment with measurable waste. For populations we show value is
frame-relative while \emph{price} is frame-independent, that the collective value throughput of a fleet is
capped by the world's entropy $\sum_a G_a\le H(X)$, and that the fleet's operating point is a Kelly portfolio
over agents selected by an emergent price. A dynamical layer gives the equations of motion and an
is/ought asymmetry --- beliefs have a target the world supplies, goals do not --- from which alignment emerges
as a control-stability condition with a closed-form residual misalignment. We then test the single-frame laws on
live language models in a \emph{pre-registered} scale-up across three task domains and a ten-model, five-family
ladder ($0.5$B--$8$B): perception mutual information tracks realized \emph{capability} rather than parameter count
(Spearman $\rho = 0.977$ pooled over 30 model$\times$domain points), out-of-sample $\Delta G$ tracks $\I(X;Y)$
(slope CI excludes $0$), and over-confidence is measurable dissipation in every domain --- the two single-frame
laws generalize. A further pre-registered test shows the out-of-sample bridge $\Delta G\sim\I(X;Y)$ is
\emph{shape-invariant}: pooled across four qualitatively different task shapes --- classification, reasoning
(GSM8K), sequential decision, and code (MBPP), all with discrete gold so $\I$ is computed, not estimated ($n=42$)
--- the slope is $0.953$, statistically indistinguishable from the classification-only value, promoting the
bridge from a demonstration toward a law. A fleet-pricing experiment is reported scoped and primary-metric-first: value-pricing recovers
cost-aware routing from first principles --- it ties good hand-tuned routing under a token budget, beats a
cost-blind router under a compute budget, and does not outperform a cost-aware engineer; its contribution is
principled measurement, not outperformance. The laws hold in the smooth, concave (diminishing-returns) regime;
threshold, satiation, and risk-seeking goals lie outside it. None of the underlying mechanisms is individually
new; the contribution is their unification under one substrate-grounded quantity, and the is/ought asymmetry that
follows.
\end{abstract}

\tableofcontents

\section{Introduction}\label{sec:intro}

Information theory began with an abstraction that looked like a loss. \citet{shannon1948} threw away
\emph{meaning} and kept only the reduction of uncertainty; from that single move came a measure ($-\log p$), a
unit (the bit), and hard limits (channel capacity, the coding theorems). The abstraction was productive
precisely because it was ruthless: by refusing to model what a message \emph{meant}, Shannon could state exactly
what any communication system can and cannot do.

We attempt the analogous move for \emph{value}. The everyday notion of value is a tangle of morality, market
price, and human preference. We throw all three away and keep one structural residue: \emph{value is the rate at
which a goal-directed agent converts a scarce physical resource into progress toward its goal.} This is
frame-relative --- it is defined only relative to a goal, the way information is defined only relative to a prior
and energy only relative to a reference frame --- but relativity does not disqualify a quantity from being
fundamental. It forces the theory to make the frame explicit. A theory of value is to agency what relativity is
to motion: the quantity is frame-dependent; the laws relating frames are universal.

The motivation is not only foundational. If, as increasingly seems likely, near-future populations of
artificial agents are each driven by an explicit objective, then a theory of value is not philosophy --- it is
the control theory for agent populations. To route, align, and govern many separated agents we need value to be
a \emph{measured} quantity with conservation laws and limits, not a verdict. This paper builds that quantity and
tests it on real agents.

\paragraph{Contributions and honesty.} We claim none of the underlying mechanisms as individually novel. The
logarithm is Bernoulli's and Kelly's; the growth-rate-as-information-rate is \citeauthor{kelly1956}'s; the
price layer is Arrow--Debreu's; the information geometry is \v{C}encov's and Amari's. In particular --- to meet
the concession here rather than leave it for a referee --- our capacity theorem (Theorem~\ref{thm:capacity}) \emph{is}
generalized Kelly; what we add is not the betting result but everything built upon it: the free-energy substrate
that supplies a unit, the cross-frame price, and the dynamics and alignment layers. Our contribution is
fourfold and we state it plainly (\S\ref{sec:related}): (i) the \emph{unification} of these accounts under one
substrate-grounded quantity; (ii) the \emph{capacity theorem} and \emph{fleet ceiling} that the unification
makes visible; (iii) the \emph{is/ought asymmetry} as a structural property of agent dynamics; and (iv) an
\emph{empirical test} on live models. Throughout, each result is followed by the assumptions that make it bite,
in Shannon's spirit of naming the model.

\bigskip
\noindent\textbf{Plan.} Part~\ref{part:statics} (statics) derives the measure, the capacity theorem, and the
Second Law for a single agent. Part~\ref{part:multi} (multi-agent) develops cross-frame price and the fleet
capacity region. Part~\ref{part:dynamics} (dynamics) gives the equations of motion and alignment-as-stability.
Part~\ref{part:experiments} (evidence) reports the synthetic and real-agent tests. Appendices collect the two
proofs.

% ======================================================================
\part{Statics: the measure, the limit, the Second Law}\label{part:statics}
% ======================================================================

\section{The measure: a logarithmic law of value}\label{sec:measure}

\paragraph{Setup.} An agent holds a divisible resource of total size $E>0$ (free energy, compute, budget,
attention --- any fungible substrate of action) and allocates it across $K$ goal-relevant channels in amounts
$e=(e_1,\dots,e_K)$, $\sum_i e_i=E$. Each channel $i$ carries a \emph{goal weight} $k_i\ge 0$ measuring how much
progress toward the goal a unit of effective work on $i$ yields. Let $V(e)$ be the value the allocation
realizes. Two independent structural arguments force the same logarithm; their agreement is the sign of a real law.

\paragraph{Confirmation I --- static: axioms uniquely fix the functional form.}
The purely static question --- what shape must $V$ take to be a consistent value measure? --- is answered
uniquely by three axioms via Cauchy's theorem (Appendix~\ref{app:log}).

\begin{axiom}[Diminishing returns]\label{ax:concave}
$V$ is increasing and strictly concave in each $e_i$: the first unit of resource on a channel is worth more than
the thousandth.
\end{axiom}
\begin{axiom}[Additivity across independent channels]\label{ax:additive}
For channels whose contributions do not interact, $V(e)=\sum_i v(e_i;k_i)$.
\end{axiom}
\begin{axiom}[Scale invariance]\label{ax:scale}
Rescaling the resource unit ($e_i\mapsto \lambda e_i$) changes total value only by an additive constant
independent of the allocation: $v(\lambda e_i;k_i)=v(e_i;k_i)+c(\lambda,k_i)$.
\end{axiom}

Axiom~\ref{ax:scale} is the load-bearing one for this route: it encodes the absence of a privileged
resource-unit (value has no absolute zero of scale, just as information has no privileged base of logarithm).
It is a Cauchy functional equation in disguise.

\paragraph{Confirmation II --- dynamic: compounding forces the same form.}
The dynamic question --- what is the correct measure of long-run gain when value compounds? --- is answered
by the ergodicity argument of \citet{peters2019} and, independently, \citet{kelly1956}: the time-average
growth of a multiplicatively-reinvested resource is the expected log of its multiplier. Section~\ref{sec:capacity}
develops this rigorously via the repeated game and horse-race structure, showing that the optimal long-run
allocation realizes $V^\star=\sum_i k_i\ln e_i$ --- the same logarithmic form, derived from dynamics with
\emph{no} scale assumption. The two confirmations address different questions: Axioms~\ref{ax:concave}--\ref{ax:scale}
fix the static functional form; the compounding route shows the dynamic optimum lands on the same shape.

\begin{theorem}[Logarithmic Value Law]\label{thm:log}
Under Axioms~\textup{\ref{ax:concave}--\ref{ax:scale}} with mild regularity, the value measure is, up to choice
of unit,
\begin{equation}
\boxed{\;V(e)\;=\;\sum_{i=1}^{K} k_i \ln e_i\;}
\end{equation}
and the value-maximizing allocation of a fixed budget $E$ is proportional to the goal weights,
$e_i^\star = E\,\hat k_i$ with $\hat k_i=k_i/\sum_j k_j$, giving
\begin{equation}
V^\star \;=\; K\ln E \;-\; K\, H(\hat k)\,,\qquad H(\hat k)=-\textstyle\sum_i \hat k_i\ln\hat k_i .
\end{equation}
\end{theorem}
\noindent Two features deserve note. First, Shannon's entropy $H$ reappears \emph{unbidden} as a penalty: a goal
spread thinly over many channels (high $H(\hat k)$) realizes less value than a focused one ---
\emph{focus is worth $K\,H(\hat k)$ nats of value}. Second, the logarithm is over-determined: the static
Axioms~\ref{ax:concave}--\ref{ax:scale} force it via the Cauchy equation (Appendix~\ref{app:log}), and
the dynamic compounding argument forces it again from \S\ref{sec:capacity} --- two routes that share no
premises and address different questions. A form over-determined this way is the signature of a real
structural law, not a modeling convenience.

\section{The capacity theorem: value is bounded by information}\label{sec:capacity}

\paragraph{The repeated game.} Now let value compound: value realized this round becomes resource next round (an
agent that achieves its goal earns more budget, trust, or compute). Each round $t$ the world settles into a
state $s\sim q$; the agent splits its budget into fractions $b=(b_1,\dots,b_n)$, $b_i\ge0$, $\sum_i b_i=1$; and
channel $i$ returns a multiplier $o_i(s)\ge 0$ per unit committed. The budget updates multiplicatively,
$E_{t+1}=E_t\sum_i b_i\,o_i(s_t)$, so $\log E_T=\log E_0+\sum_t \log\!\big(\sum_i b_i o_i(s_t)\big)$ and the
long-run \emph{value growth rate} is
\begin{equation}\label{eq:G}
G(b)\;=\;\E_{s\sim q}\!\Big[\log\textstyle\sum_i b_i\,o_i(s)\Big].
\end{equation}
That the time-average growth of a multiplicative process is the \emph{expected log} of its multiplier is the
ergodicity argument of \citet{peters2019}; it forces the logarithm again, now from dynamics rather than
psychology. $G$ is concave in $b$, so maximizing it is a well-posed convex program.

\paragraph{The horse race.} Take the canonical structure in which states are the channels, committing to channel
$i$ pays only in state $i$, and returns are quoted against a \emph{reference belief} $r$ over states,
$o_i=1/r_i$. Then $G(b)=\sum_s q_s\log(b_s/r_s)$, maximized (Gibbs) by betting one's model,
$b_s=p_s$. With a correct model $p=q$,
\begin{equation}\label{eq:Gstar}
G^\star=\sum_s q_s\log\frac{q_s}{r_s}=\KL{q}{r}.
\end{equation}
\emph{Maximal value growth equals the divergence of the agent's correct model from the baseline.} An agent whose
model is the baseline ($p=r$) grows value at rate zero: value creation is informational edge over the prevailing
expectation, nothing more.

\paragraph{Side information.} Equip the agent with a perception channel: a signal $Y$ correlated with the
world-state $X$, observed before allocating. Optimal play bets the posterior $b_s=p(s\mid y)$, with growth
$G_Y=\sum_{x,y}p(x,y)\log\frac{p(x\mid y)}{r(x)}$. The gain over acting on the prior alone is exactly mutual
information:
\begin{theorem}[Coding Theorem of Value]\label{thm:capacity}
The incremental rate at which an agent can compound value by perceiving the world through a channel $Y$ is at
most the mutual information between world-state and perception,
\begin{equation}
\boxed{\;\Delta G \;=\; G_Y-G_0 \;=\; \I(X;Y)\;,}
\end{equation}
and this bound is achieved by Bayes-proportional allocation. \textup{(Proof: Appendix~\ref{app:capacity}.)}
\end{theorem}
\noindent This has the full structure of a Shannon coding theorem. \emph{Achievability:} betting the posterior
attains $\Delta G=\I(X;Y)$. \emph{Converse:} any $b\ne p(\cdot\mid y)$ loses exactly $\KL{p(\cdot\mid y)}{b}$ per
round (Gibbs), so $\I(X;Y)$ is an upper bound. An agent cannot create value faster than it can perceive the
world; \emph{value-throughput is bounded by information-throughput.} Perception capacity is the hard ceiling on
value-generation rate --- the value analog of channel capacity.

\section{The Second Law of Value}\label{sec:secondlaw}

Drop the assumption that the agent's model is correct. Let $q$ be reality, $p$ the agent's model, $r$ the
baseline; the agent bets $b=p$ but the world is drawn from $q$. Then
\begin{equation}\label{eq:secondlaw}
\boxed{\;G_{\mathrm{actual}}=\sum_s q_s\log\frac{p_s}{r_s}
=\underbrace{\KL{q}{r}}_{\text{available potential}}-\underbrace{\KL{q}{p}}_{\text{dissipation (model error)}}\;.}
\end{equation}
Realized value is available potential minus dissipation from misalignment, every term in nats.

\emph{The name is a structural analogy, stated with explicit scope.} It holds exactly where thermodynamic
entropy does: the dissipation $\KL{q}{p}\ge 0$ is non-negative, \emph{destroyed} rather than created, and
monotonically reduced as the agent learns ($p\to q$). It \emph{breaks} where entropy does not: value here is
goal- and frame-relative --- fixed by the reference $r$ and the goal-weights $k$ --- not an observer-independent
quantity, so this is a Second Law \emph{of value}, not a corollary of statistical mechanics. We keep the name
because the formal structure --- a non-negative, non-creatable dissipation that learning minimizes --- is
genuinely shared, and bound it because the substrate is not.

Three consequences. (i) \emph{Misalignment is measurable dissipation}: confident error ($\KL{q}{p}>\KL{q}{r}$) drives
growth negative --- a wrong, certain agent actively destroys value, while calibrated humility is
value-preserving. (ii) \emph{Learning is value-recovery}: $\KL{q}{p}$ is exactly the cross-entropy excess that
predictive training minimizes, so the machine-learning objective \emph{is} the minimization of value
dissipation. (iii) \emph{Alignment is a value law}: reducing $\KL{q}{p}$ is mechanically identical to reducing
wasted value. We test \eqref{eq:secondlaw} on live models in \S\ref{sec:real} (R2), where the weakest model's
over-confidence indeed drives realized growth negative.

% ======================================================================
\part{Multi-agent: price and the fleet}\label{part:multi}
% ======================================================================

\section{Cross-frame value and the frame-independence of price}\label{sec:crossframe}

A single agent's value lives in its own frame $(q,p,r,k)$. Two agents with different goals have values that are
\emph{not} cardinally comparable --- the interpersonal-comparison impossibility of \citet{arrow1951}, which we
concede rather than circumvent. Yet they can still coordinate, exactly as economies coordinate non-comparable
utilities: through a \emph{price} on the shared resource.

\paragraph{Shadow price.} For an agent with goal mass $K=\sum_i k_i$ and budget $E$, the marginal value of
resource (the Lagrange multiplier on $\sum_i e_i=E$ in Theorem~\ref{thm:log}) is the \emph{shadow price}
$\lambda=K/E$. Two agents trade resource until their shadow prices equalize; the common $\lambda$ is a scalar
that lives in \emph{neither} agent's value frame. Hence:
\begin{quote}
\emph{Value is frame-relative; price is frame-independent.} Value is a vector in an agent's goal basis; price is
the one scalar on which separated frames agree.
\end{quote}

\paragraph{The invariant.} The transformation between two agents' value frames is a change of coordinates on the
belief simplex. The unique (up to scale) Riemannian metric invariant under such sufficient-statistic
reparametrizations is the Fisher--Rao metric \citep{cencov1982,amari2016}. It is the cross-frame invariant of
the theory: distances in belief that all agents agree on, whatever their goals. Alignment between two goals is
the cosine $\cos\theta_{ab}=\langle \hat k_a,\hat k_b\rangle$; its sign sets whether their interaction is
positive- or negative-sum.

\section{The fleet capacity region}\label{sec:fleet}

Consider $m$ agents acting on one world $X$ and drawing on one resource pool.

\paragraph{The ceiling.} Each agent's growth is bounded by its own perception (Theorem~\ref{thm:capacity}),
$G_a\le \I(X;Y_a)$. Jointly, the data-processing inequality gives the sum-rate bound
\begin{equation}\label{eq:ceiling}
\boxed{\;\sum_{a\in S} G_a \;\le\; \I(X;Y_S)\;\le\; H(X)\;}\qquad\text{for any subset }S,
\end{equation}
the collective value throughput of a fleet is capped by the entropy of the world it acts on. Headcount past
perception-saturation buys nothing. \emph{Perception diversity} (agents that perceive different slices of $X$)
lifts $\I(X;Y_S)$ toward $H(X)$; \emph{redundancy} (agents perceiving the same slice) adds exactly zero. We
confirm both on live models in \S\ref{sec:real} (R4).

\paragraph{The operating point.} Under multiplicative dynamics the shared pool multiplies each round by
$\sum_a w_a R_a$ for resource weights $w_a$, so the fleet is a \emph{Kelly portfolio over agents}: for
imperfectly-correlated agents, spreading resource and rebalancing raises the time-average growth (the
volatility-harvesting effect of \citealp{kelly1956,cover1991}). The weights that achieve the optimum are
selected by the resource \emph{price} $\pi$: price moves resource to high-shadow-price agents until $\lambda$
equalizes. Thus
\begin{quote}
the \emph{alignment matrix} $M=[\cos\theta_{ab}]$ shapes the achievable region (a cooperation dividend where
$\cos\theta>0$, a conflict tax where $\cos\theta<0$); the \emph{price} selects the operating point on it.
\end{quote}
This separates governance into two acts: \emph{shaping} the region (alignment and perception design) and
\emph{choosing} the point on it (pricing). The pricing claim is the one we test most carefully
(\S\ref{sec:real}, R5), and where the demon's precondition --- genuine diversity --- turns out to matter.

% ======================================================================
\part{Dynamics: motion and alignment}\label{part:dynamics}
% ======================================================================

\section{The equations of motion and the is/ought asymmetry}\label{sec:dynamics}

The statics fix equilibria; the dynamics say how they are approached. Three objects evolve. \emph{Beliefs} $p$
flow toward the true law $q$ by a natural-gradient (Fisher) descent on log-loss --- the realized regret of this
flow equals the cumulative dissipation $\sum_t \KL{q}{p_t}$ of \eqref{eq:secondlaw}, so \emph{learning is
value-recovery} dynamically as well as statically. \emph{Prices} $\pi$ flow toward market-clearing by
t\^atonnement (dual ascent on the resource constraint). \emph{Goals} $k$ flow under two forces: \emph{control}
(a principal pulling $k_a$ toward a target $k^\star$) and \emph{selection} (compounding reweights resource
toward higher-growth goals).

The organizing fact is an asymmetry. Beliefs and prices each flow toward a target \emph{the world supplies}:
reality provides $q$, the resource constraint provides market-clearing. Goals have \emph{no} target the world
supplies --- reality contains no fact about what ought to be valued. So beliefs and prices are \emph{learnable}
(gradient flows toward a given target); goals are only \emph{controllable} or \emph{selectable}. This is Hume's
is/ought gap recovered as a structural property of the dynamics, and \S\ref{sec:alignment} shows it is the
mathematical shape of the alignment problem.

\section{Alignment as a stability condition}\label{sec:alignment}

Let agent $a$ have goal $k_a$ and resource share $w_a$, with fleet effective goal $\bar k=\sum_a w_a k_a$.
Control pulls each goal toward target: $\dot k_a|_{\mathrm{ctrl}}=-\gamma(k_a-k^\star)$. Selection follows the
replicator $\dot w_a=w_a(G_a-\bar G)$. Summarize the environment's reward landscape over goals to first order by
its growth gradient $g:=\nabla_k G$ --- the direction in goal-space along which deviating from $k^\star$
\emph{increases} resource capture.

\begin{theorem}[Coupled-flow alignment]\label{thm:align}
By Price's equation \citep{price1970}, the selection drift of the mean is the trait--fitness covariance; with
$G_a\approx \bar G+g^\top(k_a-\bar k)$ it equals $Vg$, where $V=\Cov_a(k_a)$ is the goal-dispersion matrix.
Hence the effective goal obeys
\begin{equation}
\dot{\bar k}=-\gamma(\bar k-k^\star)+Vg,
\end{equation}
with fixed point and residual misalignment
\begin{equation}
\boxed{\;\bar k^\star=k^\star+\gamma^{-1}Vg,\qquad \|\bar k^\star-k^\star\|=\gamma^{-1}\|Vg\|\;,}
\end{equation}
and the aligned fixed point is locally stable iff $\gamma>\lambda_{\max}\!\big(\partial(Vg)/\partial\bar k\big)$.
\end{theorem}
\noindent The fleet does not in general settle at the target; it settles a distance $\|Vg\|/\gamma$ away ---
goal-dispersion times reward-pull over control gain. Perfect alignment holds iff $Vg=0$: either the environment
rewards exactly the target ($g=0$), or there is no diversity to select among ($V=0$), or control is infinite.
The governance reading is sharp:
\begin{quote}
\emph{The cheap half of alignment is incentive design, not control.} Driving $g\to0$ (aligning what \emph{pays}
with what is \emph{wanted}) removes the residual for \emph{any} $\gamma$ and \emph{any} $V$, preserving the
diversity that lifts the fleet ceiling \eqref{eq:ceiling}; raising $\gamma$ (brute-force oversight) merely
opposes a drift it never removes. Spend first on $g$.
\end{quote}
This is the actionable form of the is/ought asymmetry: because goals have no world-given target, they are
governed by what pays (selection, via $g$) and what is imposed (control, via $\gamma$).

% ======================================================================
\part{Evidence}\label{part:experiments}
% ======================================================================

\section{Synthetic validation}\label{sec:sim}

We first check that each closed-form prediction is reproduced by a Monte-Carlo world built to the theory's
assumptions. All five families pass to numerical tolerance (20/20 checks): E1, $\Delta G=\I(X;Y)$ to
$\sim\!10^{-3}$ nats; E2, the Second-Law decomposition \eqref{eq:secondlaw} exactly, including value going
negative under confident error; E3, the fleet ceiling \eqref{eq:ceiling} with diversity lifting it and
redundancy adding zero; E4, a Kelly-priced fleet beating ad-hoc allocation, including a ``Shannon's demon''
built from anti-correlated agents that individually do not grow yet collectively do; E5, cumulative dissipation
equal to regret, with a drifting world flooring dissipation at a positive value (a dynamical Second Law). These
confirm the mathematics is self-consistent and correctly derived. They are, however, circular by construction:
the worlds are drawn from the same distributions the formulas assume. The decisive test is on real agents.

\section{Real agents}\label{sec:real}

\paragraph{Instantiation.} We take a frozen, held-out 100-item decision task in which each item has a
ground-truth correct action drawn from $K=7$ classes, and an agent's output \emph{is} an action. This realizes
the perception-then-act structure of Theorem~\ref{thm:capacity} directly: the world-state $X$ is the correct
action, the perception $Y_a$ is the action model $a$ chooses. We set $r=q$ (the action marginal), so a
no-signal agent grows value at rate $0$. From each model's run we form the confusion matrix
$C_a[x,y]=\#(\text{gold}=x,\text{chosen}=y)$ and compute every quantity in nats with the same primitives used in
\S\ref{sec:sim}: $\I(X;Y_a)$ from the joint; the calibrated posterior $p_a(x\mid y)=C_a[:,y]/\sum_x C_a[:,y]$;
and realized growth $\Delta G_a=\text{mean}\,\ln\!\big(p_a(x\mid y_a)/r(x)\big)$. The agents are four local
models spanning a range of realized capability --- three of one family ($1.5$B, $3$B, $7$B) and a cross-family
model (an $8$B that is, on this task, \emph{less} capable than the $7$B) --- greedy decoding. Posteriors are
calibrated on a $48$-item fit split and scored on a $52$-item holdout; in-sample, the oracle posterior with
$r=q$ makes $\Delta G=\I$ an arithmetic identity (used only to confirm the units), so the empirical content is
the capability tracking, the out-of-sample tracking, and the fleet result.

\paragraph{R1: the bridge holds, and $\I$ tracks capability, not size.} Across the four models, mutual
information increases monotonically with realized capability (tool-accuracy), and out-of-sample value-growth
increases monotonically with $\I$ (Table~\ref{tab:r1}; with only four models we report the ordered values, not
a correlation coefficient). The decisive datum is the cross-family model D: it is \emph{larger} than C (8B vs 7B) yet
\emph{less} capable ($0.86$ vs $0.96$), and its $\I$ lands accordingly below C, near B. So the claim is not
``scale buys value'' but the sharper one: value-throughput is information-throughput, and information-throughput
is set by what the model can actually perceive.

\begin{table}[h]\centering
\begin{tabular}{lccccc}
\toprule
model & params & tool-acc & $\I(X;Y)$ & $\Delta G_{\text{in}}$ & $\Delta G_{\text{hold}}$ \\
\midrule
A & 1.5B & 0.79 & 1.277 & 1.277 & 0.915 \\
B & 3B   & 0.87 & 1.561 & 1.561 & 1.242 \\
C & 7B   & 0.96 & 1.779 & 1.779 & 1.423 \\
D (cross-family) & 8B & 0.86 & 1.513 & 1.513 & 1.109 \\
\bottomrule
\end{tabular}
\caption{R1 (nats). $\I(X;Y)$ tracks realized capability, not parameter count;
$\Delta G_{\text{hold}}$ increases with $\I$ across all four models. Model D is larger than C yet weaker, and its $\I$ lands below
C accordingly. World entropy $H(X)=1.92$ nats: C captures $93\%$, D (though larger) only $79\%$, A $66\%$.}
\label{tab:r1}
\end{table}

\paragraph{R2: over-confidence is dissipation, shrinking with capability.} Reading the same point-predictions
with a calibrated versus an over-confident posterior realizes different value; the gap is dissipation
(Table~\ref{tab:r2}). It is large for the weak models and shrinks with capability (not size: the larger-but-weaker
D dissipates more than C) --- a quantitative statement that less-capable agents must be more humble. For models A
and D, confident error drives realized growth \emph{negative}, exactly the Second Law's prediction
\eqref{eq:secondlaw}.

\begin{table}[h]\centering
\begin{tabular}{lcccc}
\toprule
model & tool-acc & $G_{\text{cal}}$ & $G_{\text{over}}$ & dissipated \\
\midrule
A (1.5B) & 0.79 & $+0.915$ & $-3.255$ & 4.17 \\
B (3B)   & 0.87 & $+1.242$ & $-0.863$ & 2.11 \\
C (7B)   & 0.96 & $+1.423$ & $+0.731$ & 0.69 \\
D (8B)   & 0.86 & $+1.109$ & $-1.660$ & 2.77 \\
\bottomrule
\end{tabular}
\caption{R2 (nats). Dissipation from false certainty falls as \emph{capability} rises (not size: the larger but
weaker D dissipates more than C); for the least-capable models over-confidence destroys value outright.}
\label{tab:r2}
\end{table}

\paragraph{R3: value per joule.} Mutual information \emph{per second of compute} falls across the family
($0.74\to0.38\to0.30$ nats/s for A,B,C), and the cross-family D is worst ($0.15$ nats/s: slowest and lower $\I$).
The weakest model is the most efficient perceiver per unit compute --- the value-theoretic reason a cheap reflex
is worth running on a narrow task. (This is the $\I$/compute curve across heterogeneous models, not a
within-model prompt ablation; see \S\ref{sec:limits}.)

\paragraph{R4: diversity beats redundancy.} Two different models jointly perceive more of $H(X)$ than either
alone; an identical greedy re-run adds exactly zero. The highest joint information comes from the strongest
diverse pair (B+C, $\I=1.869$ nats, $+0.090$ over C, closing most of the gap to $H(X)=1.921$), while the largest
\emph{lifts} come from pairing differently-wrong weak models (A+B and B+D each add ${\sim}0.2$ nats). This is the
fleet-ceiling prediction \eqref{eq:ceiling} on real agents.

\paragraph{R5: pricing beats pooling, but only where there is diversity to price (reported honestly).} On raw
out-of-sample growth, the Kelly/price fleet ($+1.328$) beats the equal-weight ensemble ($+1.315$) but
\emph{does not beat the single best model} ($+1.423$). There is no Shannon's demon here, and the theory says
why: four models on the \emph{same} task are positively-correlated agents (R4 quantifies the best residual
diversity at $+0.09$ nats), so Kelly rebalancing has no anti-correlated volatility to harvest and the best agent
dominates. This is the honest negative the synthetic E4 demon does not reproduce, because E4 was constructed with
anti-correlated agents.

The negative is confined to the cost-blind axis. Realized holdout growth \emph{per second of compute}
($\Delta G_{\text{hold}}/$s, distinct from R3's $\I/$s) inverts the ranking: model A
yields $0.531$ nats/s against C's $0.243$ (and the slow D only $0.111$), and a budget-aware price
$\propto \I_a/\text{cost}_a$ (the shadow price $\lambda=K/E$ of \S\ref{sec:crossframe}) achieves $0.328$ nats/s,
beating the best single model's density. Pricing therefore pays exactly where \S\ref{sec:fleet} says it should:
as the lever that chooses the operating point under a resource constraint, not as a free lunch that beats the
best agent when compute is unlimited. The demon needs \emph{perception diversity} --- different slices of $H(X)$,
i.e.\ specialists rather than generalists of varying strength --- which a set of same-task generalists lacks by
construction. We mark this as a falsifiable boundary, not a defeat.

\section{Generalization: a pre-registered scale-up}\label{sec:v2}

The test of \S\ref{sec:real} was one task and four models --- a demonstration, not a validation. To ask whether
the laws \emph{generalize}, we pre-registered (predictions, models, metrics, and pass/fail thresholds committed
and frozen \emph{before any model was run}) a scale-up to \textbf{three task domains} --- an intent-routing
benchmark \citep{larson2019clinc}, a multiple-choice QA benchmark \citep{hendrycks2021mmlu}, and a
topic-classification benchmark \citep{zhang2015agnews} ($K=4$--$6$; 240 items each) --- and a \textbf{ten-model ladder across five families} (Qwen, Llama, Gemma, Phi, Mistral;
$0.5$B--$8$B). Models emit a text label; calibration and all routing weights are fit on a held-out split and
scored out-of-sample, with $95\%$ bootstrap confidence intervals.

\paragraph{The headline: the bridge and the Second Law generalize.} Pooled across all 30 model$\times$domain
points, mutual information tracks realized capability at \textbf{Spearman $\rho = 0.977$} (CI $[0.916, 0.996]$),
and out-of-sample value-growth tracks mutual information with \textbf{slope $0.935$} (CI $[0.915, 0.954]$,
excluding $0$): realized $\Delta G$ out of sample \emph{is} perceived $\I(X;Y)$ (Table~\ref{tab:v2}). Per-domain
$\rho$ is $0.84$/$0.95$/$0.99$. Over-confidence dissipates value in \emph{every} domain, driving the
least-capable models sharply negative (e.g.\ $-11$ to $-14$ nats on the QA and topic tasks). The two single-frame
laws --- the coding theorem (\S\ref{sec:capacity}) and the Second Law (\S\ref{sec:secondlaw}) --- thus hold
across three task shapes and a five-family ladder, not just the one task of \S\ref{sec:real}. This is the
substantive result of the scale-up. The fleet ceiling also holds: all pairwise joint $\I(X;Y_a,Y_b) \le H(X)$,
and diverse specialists exceed redundant pairs in every domain.

\begin{table}[h]\centering
\begin{tabular}{lccc}
\toprule
quantity & pooled (30 pts) & per-domain range & threshold \\
\midrule
Spearman$(\I,\text{accuracy})$ & $0.977$ \,[$0.916,0.996$] & $0.84$--$0.99$ & $>0.8$ \;\checkmark \\
slope$(\Delta G_{\text{hold}} \sim \I)$ & $0.935$ \,[$0.915,0.954$] & --- & CI excl.\ $0$ \;\checkmark \\
over-confidence dissipation & $>0$ all domains & negative for weak models & $>0$ \;\checkmark \\
\bottomrule
\end{tabular}
\caption{Pre-registered generalization checks (95\% CIs). $I$ tracks realized capability and out-of-sample
$\Delta G$ tracks $I$, pooled across 3 domains and a 10-model, 5-family ladder.}
\label{tab:v2}
\end{table}

\paragraph{The fleet test, reported primary-metric-first and scoped.} We re-ran the pricing test
(\S\ref{sec:fleet}) in the heterogeneous, cost-constrained regime the theory favors: a \emph{specialist} fleet
(different models lead on different domains; low cross-agent error correlation) under a compute budget, routing
by value-price ($\propto \I_a/\text{cost}_a$) against round-robin, equal-weight, best-single, and a \emph{strong
hand-tuned} router that sends each query to the model most accurate on its domain.

On the \textbf{pre-registered primary cost metric (tokens)}, value-price \textbf{ties} the hand-tuned router ---
it does \emph{not} beat it (paired-bootstrap CI includes $0$). Token cost varies only $\sim\!1.2\times$ across the
ladder, so cost-aware pricing reduces to quality-first routing and selects the identical model. A token budget,
however, cannot express the real compute gradient. A \emph{post-hoc sensitivity analysis} (cache-only, not
pre-registered) re-scores on two cost proxies that do --- wall latency (hardware-dependent) and
active-params$\times$tokens ($\propto$ FLOPs $\propto$ energy, hardware-\emph{independent}) --- and value-price
beats the cost-\emph{blind} router on \emph{both} (FLOP-proxy $\Delta = +7.45$, CI $[6.33, 8.48]$): the cost-aware
advantage is \textbf{robust across cost metrics}, not an artifact of metric choice. Against a hand-tuned router
that \emph{itself} prices cost (accuracy/cost), value-price is edged out under raw latency but \textbf{exactly
ties} under the principled FLOP proxy ($\Delta = 0$: $\I/\text{FLOP}$ and accuracy/FLOP select the same model,
since $\I$ tracks accuracy) --- it \emph{matches}, and never beats, the cost-aware engineer. Against the naive
baselines (round-robin, equal-weight) value-price wins on every metric, so priced routing is no worse than ad-hoc
--- but that is the floor, not the claim.

\begin{quote}
Value-pricing derives cost-aware routing from first principles --- it ties good hand-tuned routing under a token
budget, beats a cost-blind router under a compute budget (robustly, across both latency and FLOP cost proxies),
and matches but does not outperform a cost-aware engineer. Its contribution is principled measurement and
cost-awareness, not outperformance.
\end{quote}

\noindent This is consistent with the scoping rule of \S\ref{sec:fleet}: pricing's edge is cost-awareness, which
it recovers as a law rather than a trick; it organizes and measures what good engineering already does, and we do
not claim it outperforms a competent cost-aware baseline. The generalization of the capacity and Second-Law
results --- not the routing comparison --- is what the scale-up establishes.

\subsection{Cross-shape generalization: the bridge is shape-invariant}\label{sec:shapes}

The scale-up above varied the task \emph{domain} but held the task \emph{shape} fixed: all three benchmarks are
pick-a-label classification. The sharper question is whether $\Delta G\sim\I(X;Y)$ is a property of
classification or a law that survives a change of task \emph{shape}. The coding theorem (\S\ref{sec:capacity})
predicts $\Delta G=\I(X;Y)$ for \emph{any} perception-then-act task, so the prediction is that the bridge is
shape-invariant. We pre-registered (predictions, models, metrics, thresholds frozen before any run) a cross-shape
test over three qualitatively different shapes, each chosen to have \emph{discrete ground truth} so that $\I$ is
computed exactly from a confusion matrix rather than estimated by a soft embedding proxy:
\begin{itemize}[noitemsep,topsep=2pt]
\item \textbf{reasoning} --- GSM8K \citep{cobbe2021gsm8k}: free-form chain-of-thought reduced to its integer
answer (gold mod $4$);
\item \textbf{sequential / agentic} --- a synthetic register-machine rollout: a step-by-step execution trace
reduced to its exact final state (mod $6$);
\item \textbf{code} --- MBPP \citep{austin2021mbpp}: a generated Python function, scored by
\emph{sandbox-executed} output (hash mod $4$).
\end{itemize}
The pipeline is identical to \S\ref{sec:v2} (calibration on a held-out fit split, out-of-sample scoring, $95\%$
bootstrap CIs); a six-model ladder ($0.5$B--$3$B) was run per shape --- the pre-registered $\ge\!6$ minimum, the
$7$B/$8$B extension being blocked by hardware (\S\ref{sec:limits}).

\paragraph{The bridge holds across all four shapes.} The out-of-sample slope $\Delta G_{\text{hold}}\sim\I$
passes for every new shape individually (reasoning $0.936$, sequential $1.023$, code $1.133$; each CI excludes
$0$), pooled across the three new shapes (\textbf{slope $0.956$}, CI $[0.920,1.001]$), and --- decisively ---
pooled with the three classification domains of \S\ref{sec:v2} into one relationship across \textbf{four task
shapes} ($n=42$: classification contributes the $24$ points of the frozen eight-model, three-family ladder
common to the shape runs ($8$ models $\times\,3$ domains), and the three new shapes the remaining $18$ ($6$
models $\times\,3$ shapes --- the $7$B/$8$B pair blocked on the new shapes, \S\ref{sec:limits}), so
$24+18=42$; v2's two extra families lie outside this common ladder):
Spearman$(\I,\text{accuracy})=\mathbf{0.924}$ (CI $[0.825,0.967]$) and
slope$(\Delta G_{\text{hold}}\sim\I)=\mathbf{0.953}$ (CI $[0.925,0.979]$). The cross-shape slope is
\emph{statistically indistinguishable} from the classification-only $0.935$ of \S\ref{sec:v2}: the bridge is the
\emph{same} near-unit-slope relationship whether the agent classifies, reasons, rolls out a sequential
computation, or writes code. Eleven of the twelve frozen checks pass (Table~\ref{tab:shapes}).

\begin{table}[h]\centering
\begin{tabular}{lccc}
\toprule
shape & $n$ models & Spearman$(\I,\text{acc})$ & slope$(\Delta G_{\text{hold}}\sim\I)$ \\
\midrule
classification (\S\ref{sec:v2}) & 10 & $0.977$ & $0.935$ \\
reasoning (GSM8K) & 6 & $0.943$ \;\checkmark & $0.936$ \;\checkmark \\
sequential (register-machine) & 6 & $0.886$ \;\checkmark & $1.023$ \;\checkmark \\
code (MBPP) & 6 & $0.429$ (underpowered) & $1.133$ \;\checkmark \\
\midrule
pooled new shapes & 18 & --- & $0.956$ \,[$0.920,1.001$] \\
cross-shape (4 shapes) & 42 & $0.924$ \,[$0.825,0.967$] & $0.953$ \,[$0.925,0.979$] \\
\bottomrule
\end{tabular}
\caption{Cross-shape generalization. The out-of-sample bridge $\Delta G_{\text{hold}}\sim\I$ holds across four
task shapes with a pooled slope ($0.953$) statistically indistinguishable from classification-only ($0.935$).
The classification row reports the full ten-model v2 headline; the $n=42$ cross-shape pool reuses only the
eight-model, three-family ladder common to all shapes ($24$ classification $+\,18$ new-shape points).
$11/12$ frozen checks pass; the lone miss is code's underpowered capability-ranking sub-check
(\S\ref{sec:limits}), not a shape-specific break --- the code \emph{bridge} slope itself passes.}
\label{tab:shapes}
\end{table}

\noindent This is a \emph{generalization}, not a new mechanism: the test confirms that the coding theorem's
prediction is shape-invariant, promoting the bridge from a demonstration on classification toward a law. The one
frozen check that does not clear is code's \emph{capability-ranking} sub-check
(Spearman$(\I,\text{accuracy})=0.429$), which is underpowered rather than a counterexample (\S\ref{sec:limits}
reports it in full); the code \emph{bridge} slope passes. As throughout, in-sample $\Delta G=\I$ is a
definitional identity (oracle posterior) carrying no empirical weight --- the content of this test is entirely in
the out-of-sample and cross-shape tracking.

\section{Related work}\label{sec:related}

Every component of this theory exists in some field; the contribution is the unification, the substrate
grounding, and the is/ought asymmetry. We state this explicitly and concede each component to its source.

\emph{Expected utility and diminishing value.} The axiomatic treatment of preference is \citet{vnm1944} and
\citet{savage1954}; our logarithmic measure (\S\ref{sec:measure}) coincides with \citet{bernoulli1738} and with
unit-coefficient constant-relative-risk-aversion utility, and the diminishing-returns law it encodes is
Weber--Fechner. We do not claim the form as novel; our departure is to \emph{derive} it from two independent structural
routes --- a scale-invariance Cauchy argument (static) and the compounding ergodicity of \citet{peters2019}
(dynamic) --- and to ground it in a conserved scarce resource (free energy being the natural physical
candidate) supplying a unit and a cross-frame exchange rate that expected utility lacks.

\emph{Information rate and log-optimal growth.} The closest prior art, which we credit most carefully, is
\citet{kelly1956}: the exponential growth rate of wealth is an information rate, optimized by betting one's
beliefs; \citet{breiman1961} proved asymptotic optimality and \citet{cover1991} developed universal portfolios.
Our capacity theorem (\S\ref{sec:capacity}) is generalized Kelly, and the fleet operating point (\S\ref{sec:fleet})
is a Kelly--Cover portfolio over agents. What we add is the reinterpretation of wealth as any conserved scarce resource (free energy being the physical
candidate rather than a monetary one), the cross-frame price layer, the dynamics/alignment layer, and the claim that
Kelly's theorem is the single-agent monetary \emph{special case} of a general law of value.

\emph{Reinforcement learning.} ``Value function'' is precisely defined in RL \citep{bellman1957,suttonbarto} as
expected cumulative reward; the relationship is complementary, not competitive. RL takes reward as given and
maximizes it; our goal weights are the analog of reward, but we problematize where reward comes from and how it
drifts (\S\ref{sec:dynamics}--\ref{sec:alignment}). Our theory sits beneath RL --- a candidate account of what
reward \emph{is}.

\emph{Thermodynamics of computation and the free-energy principle.} That information has thermodynamic cost is
\citet{landauer1961} and \citet{bennett1982}; \citet{jaynes1957} recast statistical mechanics as inference;
\citet{england2013} studied dissipation-driven adaptation; \citet{friston2010} casts agents as free-energy
minimizers. Our belief dynamics (\S\ref{sec:dynamics}) is free-energy-principle-like and shares the information
geometry \citep{cencov1982,amari2016}. The distinction is precise: the free-energy principle is a theory of the
perception/belief half --- the ``is'' --- whereas our value layer is the goal/value half --- the ``ought'' ---
plus a multi-agent price economics the FEP does not contain. The seam between them is the is/ought asymmetry.

\emph{General equilibrium and social choice.} Our price layer (\S\ref{sec:crossframe}) is
\citet{arrowdebreu1954,debreu1959} and its companion impossibility \citep{arrow1951}. We use these rather than
extend them: we concede that cardinal cross-agent comparison is impossible and route around it via an emergent
price. The contribution is the synthesis with the information-theoretic value layer under one substrate, and the
application to artificial-agent populations.

\emph{AI alignment.} Instrumental convergence \citep{omohundro2008,bostrom2014} we derive from the compounding
dynamics of value (\S\ref{sec:dynamics}): goal-directions correlated with resource capture are selected
regardless of terminal content. We recast alignment as a dynamical stability condition (\S\ref{sec:alignment}),
which to our knowledge is a new framing.

\section{Discussion: governing a population of agents}\label{sec:discussion}

The results compose into a control theory for populations of separated agents. \emph{Per agent}, value-generation
rate is capped by the perception mutual information $\I(X;Y_a)$ (Theorem~\ref{thm:capacity}); supplying an agent
the goal-relevant bits cheaply --- raising effective $\I$ per unit compute --- lifts its value ceiling without
raising its deliberation cost, which is the formal reason a cheap reflex can match expensive deliberation on a
narrow task (R3 measures this curve directly). \emph{Across agents}, values are not cardinally comparable, so
coordination must run through an emergent price on shared resource (\S\ref{sec:crossframe}), never a god's-eye
utility sum. \emph{Fleet design} is then two levers: cultivate perception diversity to lift the entropy ceiling
\eqref{eq:ceiling} (R4), and price resource to the best operating point (R5). The R5 result is best read not as a
weakness but as a \emph{scoping theorem} for the pricing lever: \textbf{pricing dominates ad-hoc allocation
precisely under cost constraints or imperfect agent correlation; absent both --- free compute, a single task,
and homogeneous, positively-correlated agents --- the best single agent is optimal, exactly as the theory
predicts.} A governance claim that names the regime in which it does \emph{not} win is a sharper claim, not a
softer one; and the regime in which it does win --- agents that are heterogeneous \emph{and} resource-constrained
--- is the ordinary condition of real fleets, where R3's value-density curve and R5's cost-aware price compound
into a concrete operating advantage. \emph{Alignment}, finally, is incentive design before oversight
(Theorem~\ref{thm:align}): spend first on making what pays equal what is wanted.

\section{Limits}\label{sec:limits}

We name the assumptions that bound each claim. The measure (\S\ref{sec:measure}) is the smooth, concave regime:
threshold (all-or-nothing) goals, satiation, and risk-seeking violate Axiom~\ref{ax:concave} and lie outside
the theory. The capacity theorem assumes compounding (reinvested value) and is
cleanest in the horse-race structure; general returns give a concave program without the closed-form
$\KL{q}{r}$. The fleet ceiling's exact achievable region is an open problem of network-information-theory
difficulty. The cross-frame layer concedes that cardinal interpersonal comparison is non-canonical; we do not
solve it. On the empirical side: the in-sample $\Delta G=\I$ is an arithmetic identity, not evidence ---only the
capability tracking, the out-of-sample tracking, and R5 are empirical; R3 is an $\I$/compute curve across model
scale, not a within-model prompt ablation; the calibration in R2 is constructed from the confusion matrix, not a
probability the model reported, so the dissipation measured is that of a stated belief over those predictions;
and the empirical scope is three classification domains, a ten-model five-family ladder (\S\ref{sec:v2}), and a
cross-shape extension to reasoning, sequential, and code tasks (\S\ref{sec:shapes}) --- evidence the bridge
\emph{generalizes across four task shapes}, not yet a universal law, and not independently replicated.
Degenerate runs in which a model emits no action (a serving/plumbing failure, not a decision) are detected and
excluded rather than reported as low-$\I$ capability. \emph{The one frozen cross-shape check that does not
clear is the code shape's capability-ranking sub-check} (Spearman$(\I,\text{accuracy})=0.429$, CI
$[-0.43,1.00]$): the six $\le\!3$B models cluster tightly in code-accuracy ($0.35$--$0.47$), so the rank
correlation is noise-dominated and inconclusive --- a power limitation, not a shape-specific break. The
pre-registered remedy (adding $7$B/$8$B models, far stronger coders, to widen the code capability range) could
not be run: those models stall at usable throughput on the $16$\,GB host at $512$-token chain-of-thought, the
same memory ceiling that bounds the ladder elsewhere. Crucially the code \emph{bridge} itself
(out-of-sample $\Delta G_{\text{hold}}\sim\I$) \emph{passes}; only the capability-ranking sub-check on that one
shape is underpowered, and we report it as such rather than as a counterexample or forcing the run on
inadequate hardware. The alignment theorem is mean-field with a linearized
landscape and a quasi-static dispersion $V$; a fully coupled treatment, and an endogenous reward gradient $g$ in
a closed multi-agent world, remain open.

\paragraph{What remains decisive.} The scale-up (\S\ref{sec:v2}) already ran the heterogeneous,
resource-constrained fleet test the governance claim turns on --- specialists perceiving different slices of $X$
under a compute budget --- and the answer was scoped, not triumphant: value-pricing beats naive and cost-blind
allocation but only \emph{matches} a strong cost-aware hand-tuned baseline, while collective throughput obeyed
the predicted ceiling $\sum_a G_a\le H(X)$. Two things would still be decisive, and neither is yet done. First, a
\emph{dynamic} multi-agent setting in which the goals $k$ themselves evolve over time
(\S\ref{sec:dynamics}--\ref{sec:alignment}), rather than the static classification tested here --- that is the
regime where the is/ought asymmetry and the alignment-stability theorem actually bite, and a failure of the
predicted flows there would refute the dynamical layer. Second, independent replication at larger scale. We have
shown the single-frame laws \emph{generalize} across three task shapes and a ten-model ladder; what remains is to
show the framework governs a population whose goals \emph{move}.

\section{Conclusion}\label{sec:conclusion}

Stripped of morality, price, and psychology, value is a structural quantity with a measure ($\sum_i k_i\ln e_i$),
a capacity limit ($\Delta G=\I(X;Y)$), a Second Law ($G=\KL{q}{r}-\KL{q}{p}$), a frame-independent price, a
fleet entropy ceiling, and an alignment-stability law. Shannon's own quantities --- entropy, divergence, mutual
information, the Fisher metric --- reappear unbidden throughout, which is the strongest internal evidence that
the abstraction is real rather than decorative. On live language models --- across three task domains and a
ten-model, five-family ladder, pre-registered --- the single-frame laws hold to high precision and generalize,
and the one fleet-level claim with a genuine precondition is reported with its boundary intact. The
quantity is frame-relative; the laws relating frames are universal; and they are now testable on the artificial
agents the theory was built to govern.

\appendix
\section{Proof of the Logarithmic Value Law (Theorem~\ref{thm:log})}\label{app:log}

\emph{This appendix is Confirmation~I: the static axiomatic route to $V=\sum_i k_i\ln e_i$ via the Cauchy
functional equation. Confirmation~II --- the dynamic compounding route --- is developed in
\S\ref{sec:capacity}. The two share no premises and force the same form independently.}

By Axiom~\ref{ax:additive} it suffices to find the per-channel $v(e;k)$. Axiom~\ref{ax:scale} states
$v(\lambda e;k)=v(e;k)+c(\lambda,k)$ for all $\lambda,e>0$. Fix $k$ and write $f(e):=v(e;k)$. Then
$f(\lambda e)-f(e)=c(\lambda)$ is independent of $e$. Differentiating in $\lambda$ at $\lambda=1$,
$e f'(e)=c'(1)=:k$, a separable ODE with solution $f(e)=k\ln e+\text{const}$. Concavity (Axiom~\ref{ax:concave})
fixes the sign $k\ge 0$ and rules out the additive constant's dependence on $e$. Summing over channels gives
$V(e)=\sum_i k_i\ln e_i$. Maximizing under $\sum_i e_i=E$ with multiplier $\lambda$: $\partial_{e_i}(V-\lambda
\sum_j e_j)=k_i/e_i-\lambda=0\Rightarrow e_i=k_i/\lambda$, and $\sum_i e_i=E$ gives $\lambda=K/E$ (the shadow
price) and $e_i^\star=E\hat k_i$. Substituting, $V^\star=\sum_i k_i\ln(E\hat k_i)=K\ln E+\sum_i k_i\ln\hat k_i
=K\ln E-K H(\hat k)$ after normalizing $k_i=K\hat k_i$. \hfill$\square$

\section{Proof of the Coding Theorem of Value (Theorem~\ref{thm:capacity})}\label{app:capacity}

\emph{Achievability.} With side information $Y=y$ the optimal bet is the posterior $b_s=p(s\mid y)$ (Gibbs
applied conditionally), giving conditional growth $\sum_s p(s\mid y)\log\frac{p(s\mid y)}{r(s)}$. Averaging over
$y$,
\[
G_Y=\sum_{y}p(y)\sum_s p(s\mid y)\log\frac{p(s\mid y)}{r(s)}
   =\sum_{x,y}p(x,y)\log\frac{p(x\mid y)}{r(x)}.
\]
Without side information the optimal bet is the prior $b_s=p(s)$, giving $G_0=\sum_x p(x)\log\frac{p(x)}{r(x)}
=\KL{p_X}{r}$. Subtracting,
\[
\Delta G=G_Y-G_0=\sum_{x,y}p(x,y)\log\frac{p(x\mid y)}{p(x)}=\I(X;Y).
\]
\emph{Converse.} For any allocation $b(\cdot\mid y)\ne p(\cdot\mid y)$, the conditional growth loss is
$\sum_s p(s\mid y)\log\frac{p(s\mid y)}{b(s\mid y)}=\KL{p(\cdot\mid y)}{b(\cdot\mid y)}\ge 0$ by Gibbs, with
equality iff $b=p(\cdot\mid y)$. Hence no strategy exceeds $G_0+\I(X;Y)$, so $\I(X;Y)$ is both achieved and an
upper bound on $\Delta G$. \hfill$\square$

\bibliographystyle{plainnat}
\bibliography{refs}

\end{document}
